\chapter{总结与展望}
\label{chap:conclusion}

\section{研究工作总结}

本文围绕空间公共物品博弈中的社会信息利用，研究邻居经验如何进入个体强化学习，以及信息失真后应如何调整其作用。两项工作的共同对象是动作价值更新：智能体保留独立Q表和自身试错过程，邻居信息通过附加修正改变价值判断，再影响后续行为。通过区分真实博弈、报告消息与学习更新，本文将信息可靠性与合作演化联系到可分析、可检验的具体机制上。

第一项工作在可靠报告条件下，利用邻居相对综合奖励与动作一致性修正当前已执行动作的Q值。已发表结果表明，适当引入邻居影响能够改变合作演化轨迹，在相应条件下使合作经历初期下降后恢复。扩大感知范围增加了参考来源，但其作用受到影响强度、奖励组成和状态表示共同影响。更早形成合作与更高的后期合作水平并非总是同时实现，因此需要分别观察学习过程与后期行为。

第二项工作将可靠报告这一前提放宽，明确建模持续虚报高奖励与背叛动作的消息通道。本文构造局部评分与连续门控，分别控制参考对象及正奖励优势对Q表的修正幅度。共享参数通过完整仿真的交叉熵搜索获得，执行时保持固定；个体则在新环境中从头在线学习。控制器不读取异常身份或邻居未篡改真值，并以简单过滤、合作消息优先和固定弱影响作为必要对照。

冻结后的1200次留出测试表明，完整模型在两个协同因子、干净与10\%持续失真四个条件中，相对已调固定合作规则的全程合作优势为4.105至5.255个百分点，对八项主要比较校正后的相应bootstrap区间均位于零以上。两个受扰条件的全程合作分别达到0.96774和0.98398。与此同时，固定合作评分加学习门控取得相近表现：学习评分在$r=4.4$的两个条件下增加约0.4302和0.6594个百分点，在$r=4.8$干净时反而低约0.2173个百分点，受扰时则没有得到支持完整模型更高的区间证据。由此，学习控制相对固定强度的收益与进一步学习评分的增量应分开评价。

机制检验进一步收缩了对高合作的解释。简单合作优先与学习控制器都很少选择持续报告背叛的异常来源，因而低异常选择比例并不是学习识别真伪的独有证据。冻结参数后移除显式合作指示量，两个受扰条件的全程合作均降至约0.005；移除经验残差或Q差距代理也使合作下降。这说明当前规则依赖合作语义，同时使用自身历史与价值信息。幅度校准固定规则在最终测试中的汇总绝对NI与完整模型相差约3.79\%，仍未获得相同合作水平，表明平均更新量不足以单独解释差异。

此外，1152次开发泛化运行覆盖消息类型、协同因子、异常比例、状态与规模。冻结参数能够在相同十二邻居结构下迁移到一万节点网格，但在较低协同因子、较高异常比例或四邻居感知下，完整模型可能低于简化规则。单进程计时也显示，在一万节点、相同预算下，完整模型与固定评分门控均约为每步16毫秒，高于固定合作规则约10毫秒的成本。参数量固定与局部执行支持模型迁移，却不意味着执行开销可以忽略。

\section{主要认识与研究局限}

本文的分析形成了三个相互联系的认识。首先，社会信息可以在保留标准TD步骤的同时改变价值偏好；从代数上看，附加项相当于本轮有效奖励的改变，保留TD并不自动保持原有最优策略。其次，对本文的NI形式，合作报告使附加项增加合作价值优势，背叛报告产生相反作用，非负门控负责缩放这一方向。最后，局部归一化限制单次更新幅度，并可导出有界奖励下的Q值范围；这些性质说明了更新作用与数值范围，但没有给出合作形成或策略收敛的充分条件。

实验结论首先受信息边界限制。默认状态使用未被污染的邻域声誉，攻击作用于动作与奖励报告；异常位置在单次运行内固定，发送者的真实动作仍由个体学习产生。自身动作状态与多种失真消息检验了部分边界，但状态污染、发送者协同以及根据防御反馈实时变化的攻击没有被覆盖。扩大周期网格也保持了局部拓扑，尚未检验异质度数与动态网络。

统计证据同样有明确范围。工作1以作者保存的发表图为来源，缺少逐次数据时不重建置信区间。工作2的最终测试包含30个环境重复，但只评价各类经过选型的一组冻结参数；每类三次独立优化不足以完整刻画训练失败概率。泛化批次属于开发范围探索，与最终留出结果分别报告。高低合作混合分布、未被选择的候选和低合作结果均保留，单次空间快照只用于说明局部行为。

评价指标也需要结合环境解释。标准化真实福利与合作率在本文设置中严格相等，不能作为两项独立改善证据。全程与末段指标分别描述完整过程与后期水平，策略切换和实际NI幅度提供补充信息。测试使用固定预算，后期仍存在探索和动作变化，因此本文报告有限预算内的合作表现，不将其称为已经达到理论稳态。

\section{未来研究方向}

后续研究首先可以扩展信息边界，将声誉观测的不可靠性与消息内容失真共同建模，并区分接收者能够直接观察的真实量与必须依赖他人报告的量。对于会根据防御器调整消息的发送者，需要重新定义攻击能力及训练、测试关系，再评价局部信息控制是否仍能维持合作。

另一个方向是简化控制器。固定合作评分配合学习门控已经解释大部分收益，下一步应研究哪些环境确实需要学习评分，以及较少特征是否能保留关键行为。当前特征移除检验的是冻结参数的依赖，后续可通过同预算重训比较简化结构，并在相同任务质量下测量执行成本。

最后，可以将博弈交互网络与信息感知网络分别改变，引入异质度数、动态连接和通信成本，考察群组收益结构与消息获取条件各自的影响。这类扩展应保留本文对真实交互、观测消息和价值更新的区分，使新的合作现象仍能对应清楚的机制与证据。
